\documentclass[conference]{IEEEtran}

\usepackage[utf8]{inputenc}
\usepackage{amsmath,amssymb}
\usepackage{graphicx}
\usepackage{siunitx}
\usepackage{bm}
\usepackage{hyperref}
\usepackage{cite}
\usepackage{listings}
\usepackage{xcolor}
\usepackage{placeins} 

\lstset{
    language=Matlab,
    basicstyle=\ttfamily\footnotesize,
    keywordstyle=\color{blue},
    commentstyle=\color{green!50!black},
    numbers=left,
    numberstyle=\tiny,
    breaklines=true,
    frame=single
}

\title{Terrain-Referenced UAV Navigation Using a Particle Filter and LiDAR Patch Measurements}

\author{
  \IEEEauthorblockN{Derek Yu}
  \IEEEauthorblockA{
    AAE 590 Fundamentals of Estimation - Final Project\\
    Purdue University\\
    Email: yu1481@purdue.edu
  }
}

\begin{document}
\maketitle

\begin{abstract}
This report presents a simulation study of terrain-referenced navigation for an unmanned aerial vehicle (UAV) using a particle filter and noisy LiDAR elevation measurements over a real digital elevation model (DEM) of an area in the Grand Canyon. The filter estimates the horizontal position and heading of the vehicle, treating altitude as known a priori from an external source. A square LiDAR footprint is simulated at each time step by sampling a 2D grid of ground elevations below the UAV from the DEM; these samples are corrupted by a constant bias and Gaussian noise to emulate a realistic sensor. The particle filter propagates hypotheses of the UAV state using a noisy kinematic model with uncertain speed and turn rate, and evaluates their likelihood by comparing the known terrain patch with the measured patch. Filter performance is evaluated in trajectory tracking and position error relative to the truth UAV trajectory. Results show that the particle filter is able to successfully estimate of the UAV trajectory over the terrain using the biased and noisy LiDAR data and a DEM, with position errors on the order of less than a meter (RSS) once the filter has converged.
\end{abstract}

\section{Introduction}
Terrain-referenced navigation (TRN) uses measurements of the surrounding terrain obtained using LiDAR or radar, and a stored/known digital elevation model (DEM), to estimate a vehicle's position. This approach is attractive in environments where absolute position measurements such as from GPS are degraded or unavailable. TRN has applications in low-altitude UAV navigation, cruise missile guidance, and autonomous ground vehicle localization.

In this paper, we investigate a TRN problem in which a fixed-wing UAV flies over a real terrain map in an area near a river at the Grand Canyon, represented by a GeoTIFF DEM. The vehicle carries a LiDAR-like sensor that measures ground elevation within a square footprint beneath the aircraft. The goal is to estimate the UAV's horizontal position $(x,y)$ and heading $\psi$ using a particle filter (PF), given noisy and biased LiDAR measurements and knowledge of a DEM of the surrounding area. Altitude is assumed to be known from an external sensor (e.g., barometric altimeter), so the estimation problem is restricted to planar position and heading.

A nonlinear, non-Gaussian terrain-referenced measurement model is handled using a bootstrap particle filter, and the measurement model is built directly from a high-resolution DEM read from a GeoTIFF file.\cite{USGS_GrandCanyonNP_2019} The main goals of this effort are:
\begin{itemize}
  \item Implementation of a particle filter that can successfully estimate the position of the UAV with 5m error or less relative to truth
  \item Conduct analysis of simulation runtime/compute vs performance with varying number of particles
\end{itemize}

\section{Problem Formulation}
\subsection{State and Dynamics}
The state to be estimated is the planar position and heading of the UAV:
\begin{equation}
  \bm{x}_k = 
  \begin{bmatrix}
    x_k \\ y_k \\ \psi_k
  \end{bmatrix},
\end{equation}
where $x_k$ and $y_k$ are the east and north coordinates (e.g., in meters in a Universal Transverse Mercator (UTM)-like frame), and $\psi_k$ is the heading angle (rad), measured from the $x$-axis ($0^\circ$ heading is east).

The truth trajectory is synthesized using a discrete-time kinematic model with constant nominal speed $V$ and a time-varying commanded yaw rate $\omega_{\text{cmd}}(t)$ corresponding to an ``S-turn'' maneuver:
\begin{align}
  \psi_k &= \psi_{k-1} + \omega_{\text{cmd}}(t_k)\,\Delta t, \\
  x_k &= x_{k-1} + V \cos(\psi_{k-1}) \Delta t, \\
  y_k &= y_{k-1} + V \sin(\psi_{k-1}) \Delta t,
\end{align}
with initial conditions chosen to place the trajectory within the DEM bounds. The commanded yaw profile is implemented as a sinusoid of amplitude $\omega_{\text{amp}}$ over a finite interval, followed by straight flight:
\begin{equation}
\omega_{\text{cmd}}(t) =
\begin{cases}
  \omega_{\text{amp}} \sin\left( \dfrac{2\pi t}{T_{\text{active}}} \right), & t \le T_{\text{active}},\\[4pt]
  0, & t > T_{\text{active}}.
\end{cases}
\end{equation}

In the filter's process model, the same structure is used but with additive Gaussian noise on the speed and turn rate:
\begin{align}
  v_k &= V + \delta v_k,  & \delta v_k &\sim \mathcal{N}(0,\sigma_v^2),\\
  \omega_k &= \omega_{\text{cmd}}(t_k) + \delta\omega_k, &
  \delta\omega_k &\sim \mathcal{N}(0,\sigma_\omega^2),
\end{align}

\subsection{Terrain and DEM}
The terrain is represented by a digital elevation model (DEM) loaded from a GeoTIFF file. The DEM provides ground elevation $h_{\text{DEM}}(x,y)$ in a 4.5 km by 2 km swath of the Grand Canyon; Figure~\ref{fig:dem_surface_example} shows the surface rendering, and Figure~\ref{fig:dem_pointcloud_example} shows the corresponding point cloud used later for generating measurements.

\FloatBarrier
\begin{figure}[!t]
  \centering
  \includegraphics[width=\linewidth]{figures/fig_3D_DEM_surface_example.pdf}
  \caption{3D DEM Surface}
  \label{fig:dem_surface_example}
\end{figure}

\begin{figure}[!t]
  \centering
  \includegraphics[width=\linewidth]{figures/fig_3D_DEM_pointcloud_example.pdf}
  \caption{DEM point cloud visualization}
  \label{fig:dem_pointcloud_example}
\end{figure}

\subsection{LiDAR Patch Measurement Model}
At each time step $k$, the LiDAR model essentially generates a square footprint pointcloud measurement of the terrain centered below the UAV at $(x_k,y_k)$ with half-width $L$ in each direction. A set of grid offsets is precomputed:
\begin{equation}
  \mathcal{P} = \{(\Delta x_m, \Delta y_m)\}_{m=1}^M,
\end{equation}
where $\Delta x_m,\Delta y_m \in [-L,L]$ with step size $s_{\text{patch}}$. For the truth model, the ground elevation at each offset is sampled from the DEM:
\begin{equation}
  z_{k,m}^{\text{truth}} = h_{\text{DEM}}(x_k + \Delta x_m,\ y_k + \Delta y_m).
\end{equation}

The LiDAR measurement at time $k$ is a noisy, biased version of the patch:
\begin{equation}
  z_{k,m} = z_{k,m}^{\text{truth}} + b + \eta_{k,m}, \quad
  \eta_{k,m} \sim \mathcal{N}(0,\sigma_z^2),
\end{equation}
where $b$ is a constant bias and $\sigma_z$ is the measurement noise standard deviation. A dropout probability can be applied so that entire scans are occasionally missing; in that case, the filter skips the update and only propagates the prior.

For visualization and simple diagnostics, the mean patch elevation is also computed:
\begin{equation}
  \bar{z}_k = \frac{1}{|\mathcal{V}_k|} \sum_{m \in \mathcal{V}_k} z_{k,m},
\end{equation}
where $\mathcal{V}_k$ is the set of valid points.

\section{Particle Filter Methodology}
\subsection{Bootstrap Particle Filter Structure}
A standard bootstrap particle filter is used. At each time step $k$, the algorithm maintains a set of particles
\begin{equation}
  \left\{ \bm{x}_k^{(i)}, w_k^{(i)} \right\}_{i=1}^{N_p},
\end{equation}
where $\bm{x}_k^{(i)}$ is the state hypothesis and $w_k^{(i)}$ is its associated importance weight.

The particle filter recursion consists of the following main steps:

\subsubsection{Prediction (Propagation)}
Each particle is propagated through the process model with sampled speed and turn-rate noise:
\begin{equation}
  \bm{x}_k^{(i)} \sim p(\bm{x}_k | \bm{x}_{k-1}^{(i)}, \bm{u}_k),
\end{equation}
using the kinematic equations in the previous section. This implements the importance proposal density.

\subsubsection{Measurement Prediction and Likelihood}
For each propagated particle, the predicted terrain patch is computed:
\begin{equation}
  \hat{z}_{k,m}^{(i)} = h_{\text{DEM}}(x_k^{(i)} + \Delta x_m,\ y_k^{(i)} + \Delta y_m).
\end{equation}

Given the measured patch $z_{k,m}$, the residual at each valid point is
\begin{equation}
  r_{k,m}^{(i)} = \left(z_{k,m}\right) - \hat{z}_{k,m}^{(i)}.
\end{equation}
Let $\mathcal{V}_k^{(i)}$ be the set of valid patch indices for particle $i$. The mean squared error (MSE) over the patch for that particle is
\begin{equation}
  \text{MSE}_k^{(i)} = 
  \frac{1}{|\mathcal{V}_k^{(i)}|}
  \sum_{m \in \mathcal{V}_k^{(i)}} \left(r_{k,m}^{(i)}\right)^2.
\end{equation}
Assuming independent Gaussian noise with variance $\sigma_z^2$, a Gaussian likelihood is used:
\begin{equation}
  p(\bm{z}_k|\bm{x}_k^{(i)}) \propto
  \exp\left(
    -\frac{1}{2} \frac{\text{MSE}_k^{(i)}}{\sigma_z^2}
  \right).
\end{equation}
To prevent numerical underflow, the log-likelihood is clipped below a threshold before exponentiation.

If all patch samples are invalid (e.g., due to measurement dropout), the update step is skipped and the weights are left unchanged.

\subsubsection{Weight Update and Normalization}
The unnormalized weight update follows
\begin{equation}
  \tilde{w}_k^{(i)} = w_{k-1}^{(i)} \, p(\bm{z}_k|\bm{x}_k^{(i)}),
  \label{eq:weight_update}
\end{equation}
Weights are then normalized:
\begin{equation}
  w_k^{(i)} = \frac{\tilde{w}_k^{(i)}}{\sum_{j=1}^{N_p} \tilde{w}_k^{(j)}},
  \label{eq:weight_norm}
\end{equation}

When the sum of unnormalized weights is numerically zero, indicating that all particles are effectively inconsistent with the measurement, the implementation respawns particles uniformly over the entire DEM region with uniform weights. The number of such ``respawns'' is tracked as a diagnostic.

\subsubsection{Effective Sample Size and Resampling}
Weight degeneracy is monitored using the effective sample size:
\begin{equation}
  N_{\text{eff}} = \frac{1}{\sum_{i=1}^{N_p} (w_k^{(i)})^2},
  \label{eq:neff}
\end{equation}
 When $N_{\text{eff}}$ drops below a fraction $\alpha$ of the total number of particles,
\begin{equation}
  N_{\text{eff}} < \alpha N_p,
\end{equation}
systematic resampling is performed. Systematic resampling draws $N_p$ equally spaced samples from the cumulative distribution function (CDF) of the current particle weights, resulting in a new particle set with approximately uniform weights and duplicated particles proportional to their weights.

\subsubsection{State Estimation and Covariance}
The state estimate at each time step is taken as the weighted mean of the particles:
\begin{align}
  \hat{x}_k &= \sum_{i=1}^{N_p} w_k^{(i)} x_k^{(i)}, \\
  \hat{y}_k &= \sum_{i=1}^{N_p} w_k^{(i)} y_k^{(i)}, \\
  \hat{\psi}_k &= \text{atan2}\left(
      \sum_{i} w_k^{(i)} \sin \psi_k^{(i)},
      \sum_{i} w_k^{(i)} \cos \psi_k^{(i)}
  \right),
\end{align}
where the heading is averaged on the unit circle to handle wrap-around properly.

A rough measure of position uncertainty is obtained from the weighted covariance of particle positions. The implementation computes a $1\sigma$ radial position bound:
\begin{equation}
  \sigma_{\text{pos},k} = \sqrt{
    \sum_i w_k^{(i)} (x_k^{(i)} - \hat{x}_k)^2 +
    \sum_i w_k^{(i)} (y_k^{(i)} - \hat{y}_k)^2
  }.
\end{equation}

\section{Simulation Setup}
The simulation parameters are summarized as follows:
\begin{itemize}
  \item Time step: $\Delta t = \SI{0.5}{s}$.
  \item Duration: $T = \SI{170}{s}$
  \item Speed: computed constant nominal speed $V$ selected so that the trajectory reaches a target $x$-coordinate at $T$ seconds.
  \item Yaw profile: S-turn over $T_{\text{active}} = \SI{80}{s}$ with amplitude $\omega_{\text{amp}}=0.018~\text{rad/s}$, then straight flight.
  \item Altitude: $h(t) = 500 ~\text{m}$ MSL, assumed known to the filter.
  \item LiDAR bias: $b = \SI{0.01}{m}$; noise standard deviation $\sigma_z = \SI{0.015}{m}$; dropout probability $p_{\text{drop}} = 0.02$ (occasional missing scans), consistent with small-UAV spinning LiDAR specs.\cite{velodyne_puck_datasheet}
  \item LiDAR patch: half-width $L = \SI{20}{m}$; patch step $s_{\text{patch}} = \SI{10}{m}$.
  \item Particle filter: $N_p = 10{,}000$ particles; $\sigma_v = \SI{1.5}{m/s}$; $\sigma_\omega = 0.01~\text{rad/s}$; resampling threshold $\alpha=0.5$.
\end{itemize}

The DEM is loaded from a GeoTIFF file representing a section of terrain (a Grand Canyon region). Sampling is decimated by a user-defined factor to trade off resolution and runtime.

\begin{table}[t]
  \centering
  \caption{Key simulation constants and sensor/process parameters.}
  \label{tab:constants}
  \begin{tabular}{ll}
    \hline
    Quantity & Value \\
    \hline
    Time step $\Delta t$ & $0.5~\text{s}$ \\
    Total duration $T$ & $170~\text{s}$ \\
    LiDAR patch half-width $L$ & $20~\text{m}$ \\
    Patch grid step $s_{\text{patch}}$ & $10~\text{m}$ \\
    Number of patch points $M$ & $41 \times 41$ offsets \\
    LiDAR bias $b$ & $0.01~\text{m}$ \\
    LiDAR noise $\sigma_z$ & $0.015~\text{m}$ \\
    Dropout probability $p_{\text{drop}}$ & $0.02$ \\
    Process noise $\sigma_v$ & $1.5~\text{m/s}$ \\
    Process noise $\sigma_\omega$ & $0.01~\text{rad/s}$ \\
    Resample threshold $\alpha$ & $0.5$ \\
    Nominal particles $N_p$ & $10{,}000$ (varied in sweep) \\
    DEM sample decimation & $5$  \\
    Point cloud decimation & $500$  \\
    \hline
  \end{tabular}
\end{table}

\paragraph{LiDAR Noise Model} The LiDAR noise/bias values above reflect small-UAV spinning LiDAR specifications (e.g., Velodyne Puck-class sensors report $\pm\SI{3}{cm}$ accuracy and $\pm\SI{2}{cm}$ repeatability). $\sigma_z=\SI{0.015}{m}$ or $3 cm\ 2\sigma$ captures the spec sheet LiDAR accuracy.\cite{velodyne_puck_datasheet} However, in the simulation, we will be using much larger noise and bias of 1 m and 3 m $1\sigma$ respectively because the DEM does not appear to be as accurate as this selected LiDAR model, and results in difficulty in estimating the truth trajectory.

\section{Results and Discussion}
\subsection{Trajectory Tracking}
Figure~\ref{fig:dem_traj} shows the DEM with the true and estimated trajectories overlaid. The truth trajectory starts at the west side of the swath, executes an S-turn segment, and then continues straight eastward. The particle filter can be seen struggling to estimate the trajectory initially as can be seen by the red dashed line jumping around the DEM and not tracking the black truth line. However, the filter relatively quickly converges and starts to track.

\begin{figure}[t]
  \centering
  \includegraphics[width=\linewidth]{figures/fig_dem_truth_pf.pdf}
  \caption{Terrain DEM with true UAV trajectory (solid black) and particle filter estimate (red dashed line)}
  \label{fig:dem_traj}
\end{figure}

Terrain variability is likely important for observability in TRN. It's possible that flying over the river where the elevation differences are significant and unique helped the filter to converge.

\subsection{Position Error and Uncertainty}
The position error is computed as the Euclidean distance between the estimated and true horizontal positions:
\begin{equation}
  e_k = \sqrt{(\hat{x}_k - x_k)^2 + (\hat{y}_k - y_k)^2}.
\end{equation}
Figure~\ref{fig:pos_err} plots $e_k$ versus time, along with the $3\sigma$ position uncertainty envelope $3\sigma_{\text{pos},k}$ derived from the particle spread.

\begin{figure}[t]
  \centering
  \includegraphics[width=\linewidth]{figures/fig_pos_error.pdf}
  \caption{Horizontal position error (blue line) and estimated $3\sigma$ radial uncertainty (red dashed line) versus time.}
  \label{fig:pos_err}
\end{figure}

In a correctly-tuned filter, the error should lie within the $3\sigma$ envelope for roughly $99.7\%$ of the time under Gaussian assumptions. The terrain-induced position likelihood is likely non-Gaussian and can be multi-modal depending on terrain (ie. your position/particle variance depends on your terrain, flat terrain means more variance, distinct features in terrain means less variance), so this consistency check is only approximate. Still, the comparison is still useful to look at: long stretches where $e_k$ exceeds $3\sigma_{\text{pos},k}$ could indicate underestimation of process or measurement noise, and the opposite suggests overconservative tuning.

A scalar summary of performance is the root mean square error (RMSE):
\begin{equation}
  \text{RMSE} = \sqrt{\frac{1}{N} \sum_{k=1}^{N} e_k^2 }.
\end{equation}

Results show that the overall RMSE is on the order of a few hundreds of meters, most of which is accumulated at the start where the particle filter is still attempting to converge.

\subsection{Elevation and Heading Estimation Performance}
To visualize the DEM patch measurements, Figure~\ref{fig:truth_meas} shows the LiDAR measurements generated from the truth trajectory to give a qualitative sense of how the measurement patches look.
To give an idea of what measurements are being taken by the particle filter position estimate, Figure ~\ref{fig:heading} subplot 1 shows the average elevations of the truth and PF estimate patch measurements. It can be seen that the PF measurements are not strongly tracking truth initially but eventually does converge.

Figure~\ref{fig:heading} subplot 2 shows the PF heading estimation relative to the true heading. There is a similar behavior here where initially the heading estimate is not great, but rapidly improves.

\begin{figure}[t]
  \centering
  \includegraphics[width=\linewidth]{figures/fig_truth_lidar_patches.pdf}
  \caption{Truth terrain patch measurements used for noisy LiDAR measurement generation}
  \label{fig:truth_meas}
\end{figure}

\subsection{Effect of Measurement Noise and Bias Tuning}
The small uncertainty and bias of the LiDAR model based on real products seemed to be challenging for generating a good estimate but inflating the assumed LiDAR noise $\sigma_z$ and bias $b$ improved filter convergence. With very small $\sigma_z$, it's possible the likelihood became unrealistically sharp: small DEM interpolation errors, footprint mismatch, or unmodeled attitude/altitude error were causing the particle weights to collapse, leading to frequent respawns and erratic filter estimates. Raising $\sigma_z$ widens the likelihood, tolerating the modeling mismatches and keeping $N_{\text{eff}}$ high between updates. 

\begin{figure}[t]
  \centering
  \includegraphics[width=\linewidth]{figures/fig_heading.pdf}
  \caption{True heading versus particle filter heading estimate over time}
  \label{fig:heading}
\end{figure}

\subsection{Particle Cloud Evolution}
To better understand the filter's internal behavior and do verification, snapshots of the particle cloud at selected time steps are taken. Figures~\ref{fig:particle_evolution_0}--\ref{fig:particle_evolution_136} show how the cloud initializes, de-weights bad particles and a state estimate is generated. Later frames show the cloud tracking the UAV well.

In the case of using 10,000 particles, there was exactly 1 single total resampling event. It can be seen in Figure \ref{fig:particle_evolution_0p5} that many particles have been decayed and only the ones that generate the closest patch measurements to the truth trajectory remain. You can generally see that they are all around the same elevation (green-yellow regions). As the simulation progresses, many more particles decay and we begin to estimate the truth trajectory qualitatively well.

\begin{figure}[!t]
  \centering
  \includegraphics[width=0.9\linewidth]{figures/particles_at_t=0s.png}
  \caption{Particle cloud at $t=0~\text{s}$ (Initial particle distribution).}
  \label{fig:particle_evolution_0}
\end{figure}
\vspace{-8pt}
\begin{figure}[!t]
  \centering
  \includegraphics[width=0.9\linewidth]{figures/particles_at_t=0.5s.png}
  \caption{Particle cloud at $t=0.5~\text{s}$ (many particles de-weighted).}
  \label{fig:particle_evolution_0p5}
\end{figure}
\vspace{-8pt}
\begin{figure}[!t]
  \centering
  \includegraphics[width=0.9\linewidth]{figures/particles_at_t=10s.png}
  \caption{Particle cloud at $t=10~\text{s}$ attempting to estimate UAV state}
  \label{fig:particle_evolution_10}
\end{figure}
\vspace{-8pt}
\begin{figure}[!t]
  \centering
  \includegraphics[width=0.9\linewidth]{figures/particles_at_t=136.5s.png}
  \caption{Particle cloud at $t=136.5~\text{s}$ tracking truth trajectory}
  \label{fig:particle_evolution_136}
\end{figure}

\subsection{Particle Count Performance and Runtime}
To quantify the trade between fidelity and computational cost, Table~\ref{tab:particle_sweep} summarizes runtime and accuracy for different particle counts. Higher $N_p$ reduces overall estimated error, at the cost of increased run time. However, there is a point where the root-mean-square error asymptotes while the runtime grows exponentially. The optimal particle count for balancing run time and performance appears to be at the 100,000 count for this implementation.

\begin{table}[t]
  \centering
  \caption{Particle-count sweep results.}
  \label{tab:particle_sweep}
  \begin{tabular}{lccc}
    \hline
    $N_p$ & Runtime [s] & RMSE [m] \\
    \hline
    10{,}000 & 1.02 s  & 907.7 \\
    50{,}000 & 3.76 s  & 462.1 \\
    100{,}000 & 7.14 s & 314.0 \\
    500{,}000 & 33.83 s & 320.9 \\
    \hline
  \end{tabular}
\end{table}

\section{Conclusions}
The results show that with a balanced number of particles, the filter can successfully track the UAV's trajectory over the DEM  with horizontal position errors on the order of less than a meter once converged. The position error generally is within $3\sigma$ uncertainty bounds, indicating that are able to successfully bound/predict the error.

Future extensions could include:
\begin{itemize}
  \item Doing computations using a compiled language such as C++ to reduce run time
  \item Assessing filter performance/compute time using much larger swath measurements or flying in urban environments with more distinct terrain features
  \item Comparing performance against an extended Kalman filter (EKF) or unscented Kalman filter (UKF)
\end{itemize}

\bibliographystyle{IEEEtran}
\bibliography{refs}


\clearpage
\onecolumn
\appendices
\section{Final Project MATLAB Code}

\lstinputlisting{code/final_project_particle_filter.m}
\lstinputlisting{code/viz_tif.m}

\end{document}
